\section*{अध्याय \ref{ch:b3:colloids} --- अंतरापृष्ठ, पृष्ठसक्रियक और कोलॉइड}

\begin{solution}{exo:b3:colloids:1}
$r = \qty{0.50}{\micro m}$: $\Delta p = 2 \times 0.0720/\num{5.0e-7} = \qty{2.9e5}{Pa}$; भीतर, $1.0 + 2.9 = \qty{3.9}{bar}$।
\end{solution}

\begin{solution}{exo:b3:colloids:2}
$\ln(p/p^*) = 2 \times 0.0720 \times \num{1.807e-5}/(\num{1.0e-8} \times 8.314 \times 298.15) = 0.105$; $p/p^* = 1.11$।
\end{solution}

\begin{solution}{exo:b3:colloids:3}
$\cos\theta = (40 - 20)/72 = 0.278$, $\theta = \ang{74}$: $90^\circ$ से कम, द्रव ठोस को
आंशिक रूप से \omterm{def:b3:colloids:wetting}{आर्द्र करता है}।
\end{solution}

\begin{solution}{exo:b3:colloids:4}
NaCl: $I = \qty{10}{mol.m^{-3}}$, $\kappa^{-1} = 0.96\sqrt{100/10} = \qty{3.0}{nm}$। \ce{MgSO4}: $I = \qty{40}{mol.m^{-3}}$, \qty{1.5}{nm}।
\end{solution}

\begin{solution}{exo:b3:colloids:5}
$\Gamma = \num{12.6e-3}/(2 \times 8.314 \times 298.15) = \qty{2.54e-6}{mol.m^{-2}}$; क्षेत्रफल $1/(\Gamma N_A) =
\qty{0.65}{nm^2}$ प्रति अणु।
\end{solution}

\begin{solution}{exo:b3:colloids:6}
CMC से नीचे ढलान $-12.9$ \unit{mN/m} है, $\log c$ की प्रति इकाई ($-5.0$ से $-4.0$ तक); पठार
33.4 पर है। रेखा उस तक $\log c = -4.0 + (39.1 - 33.4)/12.9 = -3.56$ पर पहुँचती है: CMC
$\approx \qty{2.8e-4}{mol/L}$।
\end{solution}

\begin{solution}{exo:b3:colloids:7}
SDS: $P = 0.350/(0.60 \times 1.67) = 0.35$, गोलों और छोटे बेलनों की
सीमा पर: CMC के निकट गोलीय \omterm{def:b3:colloids:micelle}{मिसेल}। लिपिड: $P = 0.700/(0.65 \times 1.67) = 0.64$:
द्विपरतें, अतः पुटिकाएँ।
\end{solution}

\begin{solution}{exo:b3:colloids:8}
\ce{CaCl2}: \ce{Ca^{2+}} का $60/64 = \qty{0.94}{mM}$; \ce{AlCl3}: \ce{Al^{3+}} का $60/729 = \qty{0.082}{mM}$। \ce{Na3PO4}
अपने \ce{Na+} द्वारा \omterm{def:b3:colloids:stability}{स्कंदन} करता है (लगभग \qty{20}{mM} लवण पर); फ़ॉस्फ़ेट, एक
सह-आयन, अधिशोषित होकर ऋणात्मक आवेश को बढ़ा भी सकता है।
\end{solution}

\begin{solution}{exo:b3:colloids:9}
जलरागी भाग \qty{450.0}{g/mol} में से $6 \times 44.0 + 17.0 = \qty{281.0}{g/mol}$ (पुस्तक के परमाणु भार):
HLB $= 20 \times 281.0/450.0 = 12.5$, एक तेल-में-जल पायसीकारक।
\end{solution}

\begin{solution}{exo:b3:colloids:10}
छोटे बूँदकों का तेल सतत प्रावस्था में अधिक विलेय होता है (गुणक
$\exp(2\gamma V_{\mathrm m}/rRT)$ से, जिसका घातांक \qty{50}{nm} के लिए \qty{500}{nm} की तुलना में दस गुना बड़ा है): यह
बड़े बूँदकों की ओर विसरित होता है, जो बढ़ते हैं जबकि छोटे लुप्त हो जाते हैं। यह
\omterm{def:b3:colloids:ripening}{ओस्टवाल्ड परिपक्वन} किसी \omterm{def:b3:colloids:colloid}{पायस} को स्थूल बनाता है, भले ही कोई दो बूँदक कभी आपस में न मिलें।
\end{solution}

\begin{solution}{exo:b3:colloids:11}
नवचंद्रक त्रिज्या $r/\cos\theta$ की एक गोलीय टोपी है: उसके ठीक नीचे का द्रव
$2\gamma\cos\theta/r$ कम दाब पर होता है, और स्तंभ तब तक चढ़ता है जब तक $\rho gh$ इसकी क्षतिपूर्ति न कर दे। जल के लिए
$h = 2 \times 0.072/(997 \times 9.81 \times \num{1.0e-4}) = \qty{0.15}{m}$।
\end{solution}

\begin{solution}{exo:b3:colloids:12}
\qty{100}{mM} पर \omterm{def:b3:electrode-kinetics:double-layer}{डिबाई लंबाई} \qty{1}{nm} से कम है और प्रतिकर्षण आकर्षण से पहले
समाप्त हो जाता है: कोई अवरोध नहीं बचता (\qty{50}{mM} वक्र से आगे)। मुक्त बहुलक
पृष्ठ पर कार्य नहीं करता; प्रत्यारोपित बहुलक परत, लवण चाहे जो हो, संपर्क पर
प्रतिकर्षण करती है, शृंखलाओं के अतिव्यापन पर होने वाली एन्ट्रॉपी-हानि के कारण।
\end{solution}

\begin{solution}{pb:b3:colloids:1}
\textbf{1.} क्रिस्टल जालक में प्रतिस्थापन (\ce{Al^{3+}} द्वारा \ce{Si^{4+}} का, \ce{Mg^{2+}} द्वारा \ce{Al^{3+}} का)
एक स्थायी ऋणात्मक आवेश छोड़ते हैं, जिसकी क्षतिपूर्ति विनिमेय धनायन करते हैं।
\textbf{2.} $I = \frac12(1.0 + 1.0) = \qty{1.0}{mM}$।
\textbf{3.} $\kappa^{-1} = \qty{9.6}{nm}$।
\textbf{4.} $a\kappa = 10$: द्विपरत कण की तुलना में पतली है, जो
देर्यागिन सन्निकटन की शर्त है।
\textbf{5.} $v = 2 \times 1650 \times 9.81 \times (\num{1.0e-7})^2/(9 \times \num{0.89e-3}) = \qty{4.0e-8}{m/s}$, लगभग \qty{3.5}{mm} प्रति दिन।
\textbf{6.} वान्डरवाल्स आकर्षण के प्रभावी होने से पहले ही उनकी \omterm{def:b3:electrode-kinetics:double-layer}{विसरित परतें} प्रतिकर्षण करती हैं।
\textbf{7.} $V = 2\pi\varepsilon_r\varepsilon_0a\psi^2\ln(1 + \eu^{-\kappa h}) - Aa/12h$।
\textbf{8.} लगभग $39\,kT$।
\textbf{9.} संघट्टों का $\eu^{-39} \approx 10^{-17}$ की कोटि का अंश: व्यावहारिक रूप से कभी नहीं।
\textbf{10.} यह लुप्त हो जाता है: हर संघट्ट संपर्क तक ले जाता है।
\textbf{11.} छोटी \omterm{def:b3:electrode-kinetics:double-layer}{डिबाई लंबाई} प्रतिकर्षण को छोटे अंतरालों तक सीमित कर देती है, जहाँ
$1/h$ के अनुसार बढ़ने वाला आकर्षण प्रभावी होता है।
\textbf{12.} CCC प्रति-आयन आवेश के $z^{-6}$ के अनुसार बदलती है।
\textbf{13.} $50/64 = \qty{0.78}{mM}$।
\textbf{14.} $50/729 = \qty{0.069}{mM}$।
\textbf{15.} धनायन ऋणात्मक कणों के प्रति-आयन हैं: वे
द्विपरत में संचित होकर उसका परिरक्षण करते हैं।
\textbf{16.} सल्फ़ेट एक सह-आयन है, जो कणों से प्रतिकर्षित होता है।
\textbf{17.} प्रति घन मीटर \qty{0.069}{mol} \ce{Al^{3+}}।
\textbf{18.} \qty{0.034}{mol} $\ce{Al2(SO4)3}$।
\textbf{19.} $0.0343 \times 342.3 = \qty{11.7}{g}$।
\textbf{20.} प्रति घंटे \qty{11.7}{kg}।
\textbf{21.} $\ce{Al^3+ + 3 HCO3- -> Al(OH)3 + 3 CO2}$।
\textbf{22.} \qty{1.0}{mM} में से $3 \times 0.069 = \qty{0.21}{mM}$: पाँचवाँ भाग। शेष
हाइड्रोजनकार्बोनेट और घुली हुई \ce{CO2} जल को बफ़र करते हैं: pH केवल
थोड़ा गिरता है।
\textbf{23.} अत्यधिक आवेशित ऐलुमिनियम स्पीशीज़ अधिशोषित होकर कणों का आवेश उलट देती हैं,
और कण फिर एक-दूसरे को प्रतिकर्षित करने लगते हैं।
\textbf{24.} \textbf{प्रति घन मीटर लगभग \qty{12}{g} ऐलुमिनियम सल्फ़ेट} (\qty{11.7}{g}), व्यवहार में
जल-अपघटन के लिए आवश्यक अतिरिक्त मात्रा से पहले।
\end{solution}
